\documentclass{article}
\usepackage[T1]{fontenc}
\usepackage{hyperref}

\title{OpenCode Math Copilot: A Model-Agnostic Tool Layer for Mathematical Retrieval, Computation, and Verification}
\author{[Author names and affiliations pending owner decision]}
\date{Draft protocol and manuscript skeleton}

\begin{document}
\maketitle

\begin{abstract}
This manuscript skeleton specifies an evaluation of a mathematical tool layer
combining retrieval, symbolic computation, and formal verification in an agent
workflow without changing model weights. The proposed study compares an
unassisted condition, a generic agent with comparable corpus and tool access,
and the full plugin. Its principal contrast concerns the full plugin versus
the generic agent. New experimental results are \textbf{pending}; study sizes
are proposed, not executed. Future dry-run outputs must be labeled simulated
and cannot establish mathematical performance. No improvement, universal
expertise, or security guarantee is claimed.
\end{abstract}

\section{Introduction and Scope of Claims}
The unit of evaluation is the model, agent, tools, and corpus together.
``Model-agnostic'' describes an interface not restricted to one provider,
not guaranteed compatibility or mathematical expertise for all models.
\textbf{Pending:} motivation and contributions tied to the frozen implementation
and completed evidence. The companion claim/evidence map is
\href{../research/claims-and-evidence.md}{\texttt{research/claims-and-evidence.md}};
it distinguishes historical reports, proposed controls, and pending findings.

\section{Related Work}
ToRA is an antecedent for tool-integrated mathematical reasoning with
training, rather than evidence of novelty for inference-time integration
alone~\cite{tora2023}. Lean Copilot relates to model assistance inside Lean
proof workflows~\cite{leanCopilot2024}. miniF2F is a public formal benchmark,
not a contamination-free source of private test tasks~\cite{minif2f2021}.
Lean's proof-validation reference informs the trust discussion~\cite{leanProofValidation}.
\textbf{Pending:} passage-level verification of the cited claims, a fuller
related-work review, and real study results. Metadata above is verified; no
empirical result is claimed.

\section{Architecture and Tool Surface}
\textbf{Pending:} description tied to the study's plugin commit, effective
tool inventory, corpus and index hashes, and dependency versions. Describe
retrieval coverage and excerpt limits, symbolic execution, formal checking,
and workflow integration separately. Tool availability is not correctness.

\section{Trust Model and Verification Contract}
The proposed evidence contract separates execution, compilation, target
identity, proof status under an axiom policy, and informal correspondence.
Verification requires a controlled comparison to the trusted statement and
inspection of transitive axioms, not merely process exit status or a printed
success marker. File hashes provide provenance, not semantic equivalence.
A placeholder in a challenge statement is distinct from an incomplete solution.
\textbf{Pending:} implementation-linked checks and receipts for the actual
study. The current evaluator bridge binds a fixed theorem, expected type,
imports, and axiom policy, and issues schema-2.1.0 receipts from the strict
Verifier; a six-task engineering smoke accepted all six references with no
axioms. That smoke is not a real-model result, an independent kernel check, or
an informal-correspondence review. Restricted execution and protected graders
are proposed controls, not a security proof.
Any historical Comparator report is confined to its named challenge,
commit, and environment; it does not verify all challenges or the mathematical paper.

\section{Preregistered Experimental Protocol}
\textbf{Status: proposed; completed preregistration and budget authorization pending.}
The companion protocol is
\href{../research/protocol.md}{\texttt{research/protocol.md}}.
RQ1 concerns utility, RQ2 integration value, RQ3 reliability, and RQ4
efficiency and limits. A has no tools during resolution; B is a generic agent
with the same corpus and comparable tools but no \texttt{math\_*} interfaces;
F uses the full plugin. The primary outcome is externally graded episode
success within budget; the primary contrast is F minus B.

Use family-separated development, validation, and test partitions, protected
solutions and graders, fixed Lean statements, and explicit contamination
audits. Distinguish attempts within a budgeted episode from independent
repeats. Preserve all started episodes in operational accounting and pair
conditions in a family bootstrap. Freeze prompts, budgets, rubrics, weights,
and analysis before test outcomes are available.

The proposed pilot is $30\times2\times3\times2=360$ episodes. An expanded
example is $150\times4\times3\times3=5{,}400$ episodes before ablations.
Neither has been run by this work; both require budget authorization.
The real runner is implemented with a dry-run and a fail-closed execution gate;
no real episode has been run, and a static-equivalent inventory audit is not
accepted for execution. Any dry-run manifest must be labeled \textbf{simulated},
not measured.

\section{Results Separated by Retrieval, Computation, and Proof}
\textbf{Pending. No new evaluation results are available in this skeleton.}
No simulated outputs or plausible result tables are supplied. Historical
keyword-heuristic reports are documented separately in the claim/evidence map
and must not be described as mathematical accuracy.

\subsection{Bibliographic Retrieval and Evidence}
\textbf{Pending:} source-supported grading, coverage, and failure accounting.
\subsection{Exact and Symbolic Computation}
\textbf{Pending:} correctness under registered assumptions and domains.
\subsection{Fixed-Statement Lean Proof Construction}
\textbf{Pending:} target-specific receipts; reproduction results kept separate.
\subsection{Error and Hypothesis Detection}
\textbf{Pending:} blinded rubric outcomes and adjudication records.
\subsection{Out-of-Corpus Explanation}
\textbf{Pending:} reviewed explanations and uncertainty by family/model.

\section{Ablations, Costs, and Negative Cases}
Only a leave-one-out design is proposed: remove retrieval, computation, or
resolution-time verification from F while retaining external final grading.
\textbf{Pending:} authorized ablations and cost analysis. Cost per correct
episode includes failed attempts; undefined or missing costs are not zero.
Report cold and warm caches, setup time, failures, and timeouts rather than
only successful latencies. Preserve negative results, including the historical
GPT landau-siegel keyword score of $2/3$ with versus $3/3$ without the plugin,
as a preliminary project report rather than new mathematical evidence.

\section{Limitations, Contamination, and Threats to Validity}
Public benchmarks may overlap model training or accessible proofs. Repeated
episodes are not independent tasks. Corpus access, interface differences,
human grading, unknown model revisions, small family samples, and API drift
limit interpretation. Formal validity does not establish faithful informal
translation. No claim of universal expertise or broad causal superiority follows.
\textbf{Pending:} study-specific deviations, missing data, and limitations.

\section{Reproducibility, Licenses, and Artifact Availability}
\textbf{Pending:} immutable study manifests, records, protected evaluation
release policy, receipts, dependency/image/patch hashes, license review,
redacted logs, and regeneration commands. No archive or public dataset release
is claimed. Document actual tools and AI assistance under the chosen venue's
policy; author names and affiliations remain placeholders.

\section{Conclusion Limited to Observed Evidence}
This skeleton specifies evaluation and evidence requirements. Empirical
utility, reliability gains, efficiency, and generalization remain
\textbf{pending}. A final conclusion must follow the completed evidence,
including null and negative outcomes, rather than a predetermined claim.

\bibliographystyle{plain}
\bibliography{references}
\end{document}
